\documentclass[12pt]{article}
\usepackage[T1]{fontenc}
\usepackage[utf8]{inputenc}
\usepackage{lmodern}
\usepackage{textcomp}
\usepackage{enumitem}
\usepackage{geometry}
\geometry{margin=1in}
\begin{document}
\begin{center}
Answer Key - Astronomy - Graduate Level Assessment 12 \
\small (Answers are ordered exactly as in the question paper.)
\end{center}

\section*{Multiple Choice}
\begin{enumerate}[label=\arabic*.]
\item \textbf{Answer:} D
\\ \textit{Options:} A) Neptune; B) Uranus; C) Saturn; D) all of the above
\\ \textit{Note:} Correct option: D - all of the above
\item \textbf{Answer:} A
\\ \textit{Options:} A) Hydrogen; B) Helium; C) Carbon; D) Oxygen
\\ \textit{Note:} Correct option: A - Hydrogen
\item \textbf{Answer:} D
\\ \textit{Options:} A) They are the shattered remains of collisions between planets.; B) They are chunks of rock or ice that condensed long after the planets and moons had formed.; C) They are chunks of rock or ice that were expelled from planets by volcanoes.; D) They are leftover planetesimals that never accreted into planets.
\\ \textit{Note:} Correct option: D - They are leftover planetesimals that never accreted into planets.
\item \textbf{Answer:} B
\\ \textit{Options:} A) Tunguska Siberia.; B) Chicxulub Crater Yucatan Peninsula in Mexico.; C) Quebec Canada.; D) Meteor Crater in Arizona.
\\ \textit{Note:} Correct option: B - Chicxulub Crater Yucatan Peninsula in Mexico.
\item \textbf{Answer:} C
\\ \textit{Options:} A) The nucleus contains most of the atom's mass but almost none of its volume.; B) A neutral atom always has equal numbers of electrons and protons.; C) A neutral atom always has equal numbers of neutrons and protons.; D) The electrons can only orbit at particular energy levels.
\\ \textit{Note:} Correct option: C - A neutral atom always has equal numbers of neutrons and protons.
\end{enumerate}

\section*{True / False}
\begin{enumerate}[label=\arabic*.]
\setcounter{enumi}{5}
\item \textbf{Answer:} True
\\ \textit{Options:} A) True; B) False
\\ \textit{Note:} LLM-generated true statement based on MCQ.
\item \textbf{Answer:} False
\\ \textit{Options:} A) True; B) False
\\ \textit{Note:} LLM-generated false statement. Correct answer: 'Dione'.
\item \textbf{Answer:} False
\\ \textit{Options:} A) True; B) False
\\ \textit{Note:} LLM-generated false statement. Correct answer: '9:4'.
\item \textbf{Answer:} False
\\ \textit{Options:} A) True; B) False
\\ \textit{Note:} LLM-generated false statement. Correct answer: 'A and B only'.
\item \textbf{Answer:} False
\\ \textit{Options:} A) True; B) False
\\ \textit{Note:} LLM-generated false statement. Correct answer: 'size of the planet'.
\end{enumerate}

\section*{Long Form Response}
\begin{enumerate}[label=\arabic*.]
\setcounter{enumi}{10}
\item \textbf{Answer:} Distance at which one astronomical unit measures one arcsecond from Earth.. Explain why this is the correct answer and provide supporting evidence. Reference key facts or concepts that support this choice.
\\ \textit{Note:} Long-form derived from MCQ; award credit for stating the correct idea and giving supporting reasoning.
\item \textbf{Answer:} The nuclei of comets gradually disintegrate and spread out along their orbital paths. When the Earth passes through the orbit of an comet we are bombarded by sand-sized particles that cause a meteor shower.. Explain why this is the correct answer and provide supporting evidence. Reference key facts or concepts that support this choice.
\\ \textit{Note:} Long-form derived from MCQ; award credit for stating the correct idea and giving supporting reasoning.
\end{enumerate}

\vfill
\noindent\textit{End of Answer Key}
\end{document}